\documentclass{beamer}
\usepackage{handout}
\title{}
\author{}
\date{}
\begin{document}
\fontspec{Times New Roman}
\begin{frame}
\begin{center}
\Large{Politics} \\
\vspace{3em}
\normalsize{Instructor: Tzu-Chi Hsiao} \\
\vspace{3em}
\small{Department of Political Science} \\
\vspace{1em}
\small{National Chung Cheng University}
\end{center}
\end{frame}
\begin{frame}
\begin{center}
\Large{Chapter 0 Preliminary} \\
\vspace{3em}
\normalsize{Instructor: Tzu-Chi Hsiao} \\
\vspace{3em}
\small{Department of Political Science} \\
\vspace{1em}
\small{National Chung Cheng University} \\
\end{center}
\end{frame}
\begin{frame}{}
\begin{center}
\includegraphics[width=0.5\textwidth]{instructor.png}
\end{center}
\vspace{1em}
\begin{center}
\Large{Welcome to enroll Politics !} \\
\end{center}
\end{frame}
\begin{frame}{Course Information}
\begin{itemize}
\item Instructor: Tzu-Chi Hsiao.
\item Course: Politics.
\item Credit: 3.
\item Location: R108, Social Science Building II, CCU.
\item Time: 09:10-12:00, Every Tuesday.
\item Office: R703, Social Science Building II, CCU.
\item Office Hours: 08:10-09:00, Every Tuesday.
\end{itemize}
\end{frame}
\begin{frame}{Why taking this course}
\begin{itemize}
\item Politics happens \textcolor{red}{everywhere and everyday}.
\item We interest in the origin, cause, and results.
\item We will given the fundamentals knowledge of politics.
\item Improve your skills on politics, we have examinations.
\end{itemize}
\end{frame}
\begin{frame}{Course Grading}
\begin{itemize}
\item Midterm: 50\%.
\item Final: 50\%.
\end{itemize}
\end{frame}
\begin{frame}{Course Grading (Cont'd)}
Forms of examinations:
\begin{itemize}
\item Explanation: 40\%, 5\% each.
\item Short essay: 60\%, 30\% each.
\end{itemize}
\begin{enumerate}
\item Buy a textbook and taking notes are benefit to you.
\item Download all materials to your devices and avoid hesistate to ask any problems.
\item Learning and Reviewing must cost over fourty hours per week.
\item Must avoid taking rolls, if you choose skip the lesson, must responsible for the examinations.
\end{enumerate}
\end{frame}
\begin{frame}{Examination Preparation}
\begin{itemize}
\item Memorize all the explanations.
\item Taking some notes for the short essay.
\item Browse at least two years international news close to our covered range before enter the examinations.
\end{itemize}
\setbeamercolor{block title}{bg=Green, fg=white}
\setbeamercolor{item}{fg=Green}
\begin{block}{Aboout adjusting overall score}
\begin{itemize}
\item The overall score must adjust by the instructor to achieve the Average Equility Policy (AEP).
\end{itemize}
\end{block}
\end{frame}
\begin{frame}{Course Textbook}
Main book: Politics, Newest\textsuperscript{th} edition by Andrew Heywood. \\
Reference books:
\begin{itemize}
\item Political Science: An introduction by Michael Roskin et al.
\end{itemize}
\end{frame}
\begin{frame}{Course Schedule}
\begin{center}
\begin{tabular}{|c|c|}
\hline
Week & Topic \\
\hline
1 & Chapter 00.\\
\hline
2 & Chapter 01, Chapter 02.\\
\hline
3 & Chapter 03, Chapter 04.\\
\hline
4 & Chapter 05, Chapter 06.\\
\hline
5 & Chapter 07, Chapter 08.\\
\hline
6 & Chapter 09, Chapter 10.\\
\hline
7 & \textcolor{Green}{Prepare examination}.\\
\hline
8 & \textcolor{red}{Midterm examination}, \textcolor{red}{09:10-12:10}.\\
\hline
\end{tabular}
\end{center}
\end{frame}
\begin{frame}{Course Schedule (Cont'd)}
\begin{center}
\begin{tabular}{|c|c|}
\hline
Week & Topic \\
\hline
9 & \textcolor{purple}{Midterm examination explanation}.\\
\hline
10 & Chapter 11, Chapter 12.\\
\hline
11 & Chapter 13, Chapter 14.\\
\hline
12 & Chapter 15, Chapter 16.\\
\hline
13 & Chapter 17, Chapter 18.\\
\hline
14 & Chapter 19, Chapter 20.\\
\hline
15 & \textcolor{Green}{Prepare examination}.\\
\hline
16 & \textcolor{red}{Final examination}, \textcolor{red}{09:10-12:10}.\\
\hline
\end{tabular}
\end{center}
\end{frame}
\begin{frame}{Notice of examination}
\begin{itemize}
\item Midterm examination: 09:10-12:10, Week 08.
\item Final examination: 09:10-12:10, Week 16.
\item Both examinations are held in R108, Social Science Building II, CCU.
\item Exceed 09:40 cannot enter the examination room.
\item Examinees must bring student ID card to verify.
\item Early hand in opens after 10:40.
\item \textcolor{red}{Cheating in is not allowed}.
\end{itemize}
\end{frame}
\begin{frame}{Course Policy}
\begin{enumerate}
\item Cheating or plagiarism must dropped out from the course with \textcolor{red}{zero points} and \textcolor{red}{huge demerit}. \\
\item A missed examination must gain zero points in overall score. \\
\item If you diagonsed with disabilities, please contact the instructor immediately. \\
\item The alternative way is hand in reports to replace the examinations. \\
\item Calculating the overall score by the following formula: \\
\end{enumerate}
\begin{equation*}
\text{Overall score} = \frac{\text{Midterm report} \times 0.5 + \text{Final report} \times 0.5}{2}
\end{equation*}
\end{frame}
\begin{frame}{Alternative way (Cont'd)}
\begin{enumerate}
\item Must separate some students to the corner of the examination room. \\
\item Print the short essay problems to those. \\  
\item Their tasks are typing the answers in their personnal computer. \\
\item With Biaukai font in 12pt size in Mandarin, and Times New Roman in 12pt size in English or Numeric. \\
\item Convert to a PDF file and send to the instructor's email before 12:00. \\
\item Those who use this way must excluded to the Average Equility Policy (AEP). \\
\end{enumerate}
\end{frame}
\begin{frame}{Sharing time}
\begin{itemize}
\item What do we need to study in University.
\item What do we need to study in Political Science.
\end{itemize}
\end{frame}
\begin{frame}{Sharing time}
\begin{itemize}
\item What do we need to study in University.
\item \textcolor{gray}{What do we need to study in Political Science.}
\end{itemize}
\end{frame}
\begin{frame}{Study in University}
\begin{center}
\includegraphics[width=0.25\textwidth]{target.png}
\end{center}
\begin{itemize}
\item \textcolor{red}{Target} is the most important thing in your life.
\item We need to find how we can achieve our target via University.
\item Laid flat is a good idea ? It seems not.
\item An university is a place to pursue dream, knowledge, and truth.
\item Sounds to idealize, but it is realize in our life. 
\end{itemize}
\end{frame}
\begin{frame}{Sharing time}
\begin{itemize}
\item \textcolor{gray}{What do we need to study in University.}
\item What do we need to study in Political Science.
\end{itemize}
\end{frame}
\begin{frame}{Study in Political Science}
\begin{center}
\includegraphics[width=0.25\textwidth]{motivate.png}
\end{center}
\begin{itemize}
\item Political Science alikes a \textcolor{red}{jumper} department.
\item We learned in width but less in depth.
\item Here can give you a wide of thinking and arguing skills.
\item Training a political way thinker to discover our world.
\item Upper dimension thinking helps us to use the lower dimension attacking.
\end{itemize}
\end{frame}
\begin{frame}{}
\begin{center}
\Large{Feel free to contact me via \href{mailto:politics.tchsiao@gmail.com}{here}.}
\end{center}
\end{frame}
\begin{frame}{}
\begin{center}
\Large{End of Chapter 0}
\end{center}
\end{frame}
\end{document}
